% GENERATED FILE. Edit canonical lessons, then run npm run generate:books.
% canonical-source-hash: fnv1a64:ef0985b0d7f58a52
% canonical-lessons: LA-C273-nuptiae, LA-C273-coniunx, LA-C273-parens, LA-C273-nepos, LA-C273-neptis

\chapter{Wedding, Husband or wife, Parent, Grandson, Granddaughter}
\label{ch:la-wedding-husband-or-wife-parent-grandson-granddaughter}
\hlchaptermodality{273}

\begin{chapteropening}
\textbf{By the end of this chapter:} \emph{I can say a wedding, a husband or wife, a parent, a grandson and a granddaughter.}

Complete the last lesson of chapter 273: I can say a wedding, a husband or wife, a parent, a grandson and a granddaughter.
\end{chapteropening}

\section[nūptiae, nūptiārum]{nūptiae, nūptiārum — a wedding}

\label{lesson:LA-C273-nuptiae}



\begin{quote}
Before the new one: say the Latin for an envoy, then the Latin for a clerk.
\end{quote}

\subsection*{You'll want to know: nūptiae, nūptiārum}
\textbf{nūptiae, nūptiārum} — \textquotedblleft{}a wedding\textquotedblright{}.

\begin{cousinweb}
\textbf{nūptiae} is built on \emph{nūpta}, \textquotedblleft{}a bride\textquotedblright{} (Lewis \& Short), the participle of \emph{nūbō}, \textquotedblleft{}to marry\textquotedblright{} (EtymDB 2.1). English \emph{nuptial} comes from it.
\end{cousinweb}

\begin{grammarlens}[title={What you've built}]
One more word for this chapter.
\end{grammarlens}

\subsection*{Your turn}
\begin{itemize}
\raggedright
  \item \emph{Say it:} \emph{nūptiae, nūptiārum}
  \item \emph{Say it:} \emph{nūptiae, nūptiārum}, once more
  \item \emph{Say it:} say \emph{nūptiae, nūptiārum}
  \item \emph{Recall:} say the Latin for an envoy, then the Latin for a clerk, then say \emph{nūptiae, nūptiārum} again
\end{itemize}

\begin{tcolorbox}[breakable,colback=teal!4,colframe=teal!35!black,title={Before you move on}]
What does nūptiae, nūptiārum mean? (\textquotedblleft{}A wedding\textquotedblright{}.) Say it once more.
\end{tcolorbox}
\section[coniūnx, coniugis]{coniūnx, coniugis — a husband or wife, a spouse}

\label{lesson:LA-C273-coniunx}



\begin{quote}
Before the new one: say the Latin for a clerk, then the Latin for a wedding.
\end{quote}

\subsection*{You'll want to know: coniūnx, coniugis}
\textbf{coniūnx, coniugis} — \textquotedblleft{}a husband or wife, a spouse\textquotedblright{}.

\begin{cousinweb}
\textbf{coniūnx} is built on \emph{coniungō}, \textquotedblleft{}to join together\textquotedblright{}, from \emph{con-}, \textquotedblleft{}together\textquotedblright{}, and \emph{iungō}, \textquotedblleft{}to join\textquotedblright{}: one joined in marriage. Its root is the one in \emph{iugum}, \textquotedblleft{}a yoke\textquotedblright{}.
\end{cousinweb}

\begin{grammarlens}[title={What you've built}]
One more word for this chapter.
\end{grammarlens}

\subsection*{Your turn}
\begin{itemize}
\raggedright
  \item \emph{Say it:} \emph{coniūnx, coniugis}
  \item \emph{Say it:} \emph{coniūnx, coniugis}, once more
  \item \emph{Say it:} say \emph{coniūnx, coniugis}
  \item \emph{Recall:} say the Latin for a clerk, then the Latin for a wedding, then say \emph{coniūnx, coniugis} again
\end{itemize}

\begin{tcolorbox}[breakable,colback=teal!4,colframe=teal!35!black,title={Before you move on}]
What does coniūnx, coniugis mean? (\textquotedblleft{}A husband or wife, a spouse\textquotedblright{}.) Say it once more.
\end{tcolorbox}
\section[parēns, parentis]{parēns, parentis — a parent}

\label{lesson:LA-C273-parens}



\begin{quote}
Before the new one: say the Latin for a wedding, then the Latin for a husband or wife.
\end{quote}

\subsection*{You'll want to know: parēns, parentis}
\textbf{parēns, parentis} — \textquotedblleft{}a parent\textquotedblright{}.

\begin{cousinweb}
\textbf{parēns} is the present participle of \emph{pariō}, \textquotedblleft{}to bring forth, to give birth to\textquotedblright{} (EtymDB 2.1): the one who brings forth. It is a different word from \emph{pārēns}, \textquotedblleft{}obeying\textquotedblright{}, from \emph{pāreō}, \textquotedblleft{}to obey\textquotedblright{}.
\end{cousinweb}

\begin{grammarlens}[title={What you've built}]
One more word for this chapter.
\end{grammarlens}

\subsection*{Your turn}
\begin{itemize}
\raggedright
  \item \emph{Say it:} \emph{parēns, parentis}
  \item \emph{Say it:} \emph{parēns, parentis}, once more
  \item \emph{Say it:} say \emph{parēns, parentis}
  \item \emph{Recall:} say the Latin for a wedding, then the Latin for a husband or wife, then say \emph{parēns, parentis} again
\end{itemize}

\begin{tcolorbox}[breakable,colback=teal!4,colframe=teal!35!black,title={Before you move on}]
What does parēns, parentis mean? (\textquotedblleft{}A parent\textquotedblright{}.) Say it once more.
\end{tcolorbox}
\section[nepōs, nepōtis]{nepōs, nepōtis — a grandson}

\label{lesson:LA-C273-nepos}



\begin{quote}
Before the new one: say the Latin for a husband or wife, then the Latin for a parent.
\end{quote}

\subsection*{You'll want to know: nepōs, nepōtis}
\textbf{nepōs, nepōtis} — \textquotedblleft{}a grandson\textquotedblright{}.

\begin{cousinweb}
\textbf{nepōs} is inherited, usually traced to Proto-Indo-European \emph{*nepōts}, \textquotedblleft{}grandson, nephew\textquotedblright{} (EtymDB 2.1); German \emph{Neffe}, \textquotedblleft{}nephew\textquotedblright{}, is a cousin (Lewis \& Short). English \emph{nephew} comes from \emph{nepōs} through French.
\end{cousinweb}

\begin{grammarlens}[title={What you've built}]
One more word for this chapter.
\end{grammarlens}

\subsection*{Your turn}
\begin{itemize}
\raggedright
  \item \emph{Say it:} \emph{nepōs, nepōtis}
  \item \emph{Say it:} \emph{nepōs, nepōtis}, once more
  \item \emph{Say it:} say \emph{nepōs, nepōtis}
  \item \emph{Recall:} say the Latin for a husband or wife, then the Latin for a parent, then say \emph{nepōs, nepōtis} again
\end{itemize}

\begin{tcolorbox}[breakable,colback=teal!4,colframe=teal!35!black,title={Before you move on}]
What does nepōs, nepōtis mean? (\textquotedblleft{}A grandson\textquotedblright{}.) Say it once more.
\end{tcolorbox}
\section[neptis, neptis]{neptis, neptis — a granddaughter}

\label{lesson:LA-C273-neptis}



\begin{quote}
Before the new one: say the Latin for a parent, then the Latin for a grandson.
\end{quote}

\subsection*{You'll want to know: neptis, neptis}
\textbf{neptis, neptis} — \textquotedblleft{}a granddaughter\textquotedblright{}.

\begin{cousinweb}
\textbf{neptis} is inherited, usually traced to Proto-Indo-European \emph{*neptih2}, \textquotedblleft{}granddaughter, niece\textquotedblright{} (EtymDB 2.1), the feminine partner of \emph{nepōs}, \textquotedblleft{}grandson\textquotedblright{} (Lewis \& Short). English \emph{niece} comes from it through French.
\end{cousinweb}

\begin{grammarlens}[title={What you've built}]
One more word for this chapter.
\end{grammarlens}

\subsection*{Your turn}
\begin{itemize}
\raggedright
  \item \emph{Say it:} \emph{neptis, neptis}
  \item \emph{Say it:} \emph{neptis, neptis}, once more
  \item \emph{Say it:} say \emph{neptis, neptis}
  \item \emph{Recall:} say the Latin for a parent, then the Latin for a grandson, then say \emph{neptis, neptis} again
\end{itemize}

\begin{tcolorbox}[breakable,colback=teal!4,colframe=teal!35!black,title={Before you move on}]
What does neptis, neptis mean? (\textquotedblleft{}A granddaughter\textquotedblright{}.) Say it once more.
\end{tcolorbox}
